\documentclass[11pt]{article}
\usepackage[margin=1in]{geometry}
\usepackage[T1]{fontenc}
\usepackage{longtable}
\usepackage{array}
\begin{document}
\section*{FLCR-01 - Grounding Contract And Axiom Discipline}

**Series:** Formalizing LCR, Unifying Standard Models
**Artifact:** formal paper source
**Status:** first-pass rich rewrite; requires final citation and build pass.
**Claim posture:** maximal local claim posture with explicit lane boundaries.

\subsection*{Abstract}

This paper establishes the publication contract for the FLCR series. The core contribution is a conservative-extension discipline: imported mathematics remains external mathematics, while LCR contributes only declared bridges, receipt-bearing finite constructions, and typed claim lanes. The result is a burden system that allows maximal claims to be stated without silently changing their evidentiary type.

\subsection*{Keywords}

conservative extension; claim typing; receipts; finite proof; burden contract.

\subsection*{External Reader Orientation}

An outside reviewer should read this paper as a formal-system construction paper. Its local object is **Grounding Contract And Axiom Discipline**. The paper's immediate contribution is: Establishes the burden contract, only-addition rule, and claim-lane discipline.

The nearest external anchors are finite-state reasoning, typed interfaces, proof-carrying records, conservative extension, and ordinary mathematical citation discipline. LCR terms are therefore internal labels for the construction unless a row explicitly imports an external theorem, citation, dataset, or calibration target.

It does not ask the reader to accept any Standard Model or literal-physics identification at this layer. Such mappings are deferred to the translation papers and must carry their own evidence.

\subsection*{Reviewer Compression}

**Core object.** Grounding Contract And Axiom Discipline.

**Primary result.** This paper contributes the local formal object and its safe downstream-use rule.

**Primary non-result.** This paper does not claim direct identity with Standard Model or physical objects.

**Review strategy.** Evaluate the formal claim rows first, then inspect the receipt/citation/calibration rows that support each claim. Appendices provide routing and implementation context; they should not be read as stronger than their lane permits.

\subsection*{Peer Reviewer Assessment}

Reviewer assumption: the reviewer has been briefed on the FLCR structure, claim lanes, CAM/crystal routing, forge/action surfaces, and the distinction between internal formal results and external physical calibration.

**Recommendation.** major revision, technically promising.

**Review score vector.** orientation=2, claim\_granularity=2, evidence\_visibility=2, falsifier\_visibility=2, journal\_apparatus=1, base\_layer\_boundary=2.

**Formal claim rows reviewed.** 3.

\subsubsection*{Reviewer Strength Finding}

The manuscript is now readable as a bounded formal-system paper: it identifies its local object, separates internal construction from physical interpretation, and exposes claim lanes for review.

\subsubsection*{Major Revision Requests}

\noindent{}-- Convert the strongest proof sketch into numbered definitions, lemmas, and propositions with explicit dependencies.\\
\noindent{}-- Replace placeholder source classes with final citation keys, receipt hashes, or dataset identifiers wherever possible.\\
\noindent{}-- Move repeated pass-derived material into a compact evidence appendix and keep the main body focused on the paper-local argument.\\
\noindent{}-- Add at least one worked example showing the local object, legal operation, receipt, residue, and downstream import rule.\\

\subsubsection*{Minor Revision Requests}

\noindent{}-- Normalize theorem capitalization and punctuation.\\
\noindent{}-- Add final figure/table captions where the text currently describes diagrams schematically.\\
\noindent{}-- Ensure every acronym and local term appears in the paper-local glossary.\\

\subsubsection*{Acceptance Blockers}

\noindent{}-- A reviewer can accept the formal contribution without accepting later physics claims, but only if final citations and receipt identifiers are attached.\\
\noindent{}-- The proof sketch must be promoted into a formal proof or explicitly labeled as a construction argument.\\


\subsection*{1. Contribution And Scope}

\noindent{}-- Defines the FLCR burden contract and makes claim lanes mandatory.\\
\noindent{}-- Separates imported theorem, internal finite proof, CAM/crystal reapplication, calibration, and falsifier lanes.\\
\noindent{}-- Treats the paper corpus as proof-carrying text: claim, evidence, boundary, and falsifier travel together.\\
\noindent{}-- Installs the rule that finite exhaustive enumeration is proof only for its stated finite domain.\\

\noindent{}-- Defines open obligations as active research channels rather than restrictive\\
  labels.
\noindent{}-- Defines analog workbooks as universal STEM access kits as well as proof aids.\\
\noindent{}-- States that guardrails control AI automation and claim promotion, not human\\
  creativity or hypothesis generation.

This paper belongs to the LCR-native construction band. Physics and Standard Model language, when it appears, is only a deferred mapping cue, analogy, or explicitly bounded calibration target. The paper's base object must remain stable enough for later papers to cite without importing a stronger physical claim by implication.

\subsection*{2. Source Routing}

Primary routed sources:

\noindent{}-- `00\_Established\_Grounding\_and\_Axiom\_Contract.md`\\
\noindent{}-- `supplements/00\_INTERNAL\_REASONING\_SUPPLEMENT.md`\\

Cross-corpus evidence classes:

\noindent{}-- `CAM\_CLAIM\_100\_SCORE\_AUDIT.md`\\
\noindent{}-- `NSIT\_TOOL\_CLOSURE\_MATRIX.md`\\
\noindent{}-- `NSIT\_REASONED\_CLOSURE\_PASS.md`\\
\noindent{}-- `PAPER\_REWORKS\_CRYSTAL\_PROJECTION.json`\\
\noindent{}-- standards, glue, window, and node databases where applicable\\
\noindent{}-- notebook-derived review prompts and orbital source manifests\\

\subsection*{3. Definitions}

\noindent{}-- **Conservative extension contract.** A rule that imported theories are not modified by FLCR prose; new structure is admitted only as declared bridges and receipt-bound constructions.\\
\noindent{}-- **Claim lane.** A type attached to each claim that states what kind of evidence may consume or promote it.\\
\noindent{}-- **Receipt.** A replayable record containing input, operation, output, residue, boundary, and ideally a content hash.\\
\noindent{}-- **Open obligation.** A named next need or unresolved residue, not a truth-value label and not a silent rejection.\\
\noindent{}-- **Analog access kit.** A workbook protocol designed so the underlying STEM\\
  object can be reconstructed through available media: physical arrangement,
  gesture, counting, speech, drawing, memory, or tool-assisted replay.
\noindent{}-- **Automation guardrail.** A rule for agents, validators, scripts, and\\
  publication workflows that prevents silent promotion across evidence lanes.
  It does not limit human hypothesis formation.
\noindent{}-- **Prestate.** The declared current state of evidence and structure that acts\\
  as the first principle for the next exploration.

\subsection*{4. Formal Claims}

\noindent\texttt{| Claim | Lane | Statement |}\\
\noindent\texttt{|-------|------|-----------|}\\
\noindent\texttt{| Theorem 1.1 | `normal\_form\_result` | The FLCR corpus is admissible only as a typed extension stack: every promoted claim must declare its lane, source, evidence, and boundary. |}\\
\noindent\texttt{| Proposition 1.2 | `standard\_theorem\_citation\_bound\_result` | For finite domains such as the eight binary LCR states, exhaustive enumeration is complete proof of the enumerated property. |}\\
\noindent\texttt{| Protocol 1.3 | `falsifier\_or\_open\_obligation` | A claim without a receipt, citation, calibration datum, CAM route, or falsifier boundary remains a staged obligation. |}\\

\subsection*{Reviewer Claim Ledger}

This ledger expands the formal-claims table into review actions. It is intended to make proof granularity explicit: what is being claimed, what evidence type can support it, what boundary remains, and what the next review action is.

\noindent\texttt{| Formal row | Type | Claim in review terms | Evidence required | Boundary | Next review action |}\\
\noindent\texttt{|---|---|---|---|---|---|}\\
\noindent\texttt{| Theorem 1.1 | formal construction | The FLCR corpus is admissible only as a typed extension stack: every promoted claim must declare its lane, source, evidence, and boundary | definitions, normal-form construction, con}\\
\noindent\texttt{| Proposition 1.2 | formal construction | For finite domains such as the eight binary LCR states, exhaustive enumeration is complete proof of the enumerated property | named external theorem, definition layer, citation t}\\
\noindent\texttt{| Protocol 1.3 | obligation/falsifier | A claim without a receipt, citation, calibration datum, CAM route, or falsifier boundary remains a staged obligation | explicit missing derivation, adapter, receipt, dataset, compa}\\

\subsection*{Claim Granularity And Review Table}

The table below separates the claim types so the paper is reviewable without accepting the whole FLCR vocabulary at once.

\noindent\texttt{| Claim type | What this paper may claim | Acceptance test | What is not claimed by that row |}\\
\noindent\texttt{|---|---|---|---|}\\
\noindent\texttt{| Formal-system result | `FLCR-01` defines or uses **Grounding Contract And Axiom Discipline** as a typed FLCR object with declared inputs, operations, outputs, and residue. | Definitions, formal claims, construction ste}\\
\noindent\texttt{| Computational or receipt-bound result | Enumerations, TarPit runs, CAM/crystal routes, verifier rows, and generated artifacts are claims about the stated finite or executable domain. | A named receipt, validator, manif}\\
\noindent\texttt{| Imported theorem or external definition | Imported mathematics remains an external theorem/citation target; this paper may use it only through declared mappings and conservative-extension discipline. | The source theor}\\
\noindent\texttt{| Interpretive bridge or analogy | Analog, workbook, visual, or translation language may explain why the construction is useful. | The analogy preserves the relevant structure and does not promote itself into a theorem. }\\
\noindent\texttt{| Physical or calibration-facing claim | Physical or Standard Model identity is not claimed here unless the paper contains an explicit calibration row; the default status is deferred mapping. | A dataset, citation, calib}\\
\noindent\texttt{| Open obligation or falsifier | A missing derivation, adapter, receipt, dataset, or failed comparator is a named research channel. | The next binding action and the reason the stronger claim is not closed are stated. | }\\

\subsection*{5. Paper-Specific Development}

1. Start with imported mathematics as fixed external burden: FLCR may cite it, but may not silently mutate it.
2. Declare the new local contribution as bridge plus receipt discipline, not replacement theory.
3. Require each downstream claim to carry a lane, evidence input, boundary, and falsifier.
4. Allow later papers to strengthen claims only by adding evidence, never by erasing the original boundary.

\subsubsection*{Proof Sketch}

The argument is a governance proof. If every claim has a type and every consumer is required to check that type, then claims cannot be promoted across evidentiary boundaries by prose alone. This is the same structural move used in typed interfaces and conservative extensions: the extension may add named objects, but it cannot make imported statements mean new things without an explicit bridge.

\subsubsection*{Research Posture And Intended Direction}

This contract is not a timidity rule. The series may and should state the
largest coherent version of its intended claims. The formal layer displays the
parts already proved or source-bound; the obligation layer names the channels
between those proved parts and the larger claim. The obligation is therefore a
research route, not a prohibition.

The analog toolkit is part of this posture. It exists to let the same state be
handled without dependence on a specific technology stack. A valid workbook
should eventually be translatable into field exercises that work without
electricity, internet, computers, or paper, while still preserving the formal
state being addressed.

The physics direction is explicit. Early LCR-native papers may avoid premature
physical identification at their local level, but the later translation and
synthesis papers are expected to make the machinery increasingly literal to
physics through adapters, standard formalism, calibration, and falsifiers.

\subsection*{6. Construction}

The construction is intentionally staged. First, the paper names the finite or
typed object that can be inspected directly. Second, it declares the admissible
operations over that object. Third, it records the receipt or source class that
allows another paper to consume the result. Fourth, it names the residue that
must not be silently promoted.

For this paper, the operative object is `Grounding Contract And Axiom Discipline`. Its safe consumption
rule is:

\begin{verbatim}
source object -> declared operation -> receipt/lane -> downstream import
\end{verbatim}

The unsafe consumption rule is:

\begin{verbatim}
source phrase -> downstream identity claim
\end{verbatim}

That second pattern is explicitly rejected by FLCR-01 and must be rewritten as
a claim-lane promotion with evidence.

\subsection*{7. Evidence And Receipts}

\noindent\texttt{| Evidence source | What it contributes |}\\
\noindent\texttt{|-----------------|---------------------|}\\
\noindent\texttt{| Paper 00 legacy body | Grounding contract, imported-theorem burden, only-addition rule. |}\\
\noindent\texttt{| Internal supplement 00 | Conservative-extension reading, proof-carrying-code framing, finite-model-checking note. |}\\
\noindent\texttt{| CAM Claim 100 Score Audit | Paper 00 scored 96, with remaining work limited to citation polish and burden table finalization. |}\\
\noindent\texttt{| NSIT tool matrix | Standards-level validation lanes for claim envelopes and content-addressed receipts. |}\\

\subsection*{Appendix A. Recursive Capability Reapplication Review}

This review treats the newly admitted capability families as operators. The question is not only what each source says, but what it solves when the source is recursively reapplied through CAM, claim lanes, forge admission, receipt loops, and paper-to-paper routing.

\subsubsection*{Operator Effects}

\noindent\texttt{| Operator | Standalone result | Recursive result | Newly resolved state | Remaining boundary |}\\
\noindent\texttt{|---|---|---|---|---|}\\
\noindent\texttt{| REAPPLY-PROOFVALIDATED-PAPER-ASSEMBLY | The assembly protocol defines a paper as a formal/tool/workbook/verifier unit rather than a prose-only artifact. | When reapplied, the publication series becomes self-validating:}\\

The practical consequence is that several prior gaps are now better classified. Some are closed as routing, admission, finite-address, or executable-publication problems. The residues that remain are sharper: exhaustive source intake, physical calibration, full paper-local runner parity, external measurement, and domain-specific validation.

\subsection*{Appendix B. Broad Capability Intake}

This addendum admits non-obvious source material by function rather than filename. Each row states what the source now lets the paper do, while detailed receipt provenance remains in the evidence report instead of the formal claim body.

\noindent\texttt{| Capability | Lane | Paper effect | Boundary |}\\
\noindent\texttt{|---|---|---|---|}\\
\noindent\texttt{| CAP-PROOFVALIDATED-PAPER-ASSEMBLY | `receipt\_bound\_internal\_result` | Adds an older but still useful production contract for papers as executable triads: formal block, tool binding, runnable verifier, workbook analogue}\\

The resulting claim posture is broader but still auditable: a paper can import a capability when the evidence lane is satisfied, and any remaining gap is named as an adapter, calibration, rerun, or falsifier obligation.

\subsection*{Appendix C. Implementation Intake}

This addendum binds implementation work that was not fully represented in the earlier paper routing. The rows below are not pasted source text; they are evidence-lane imports that change what the paper can claim now.

\noindent\texttt{| Implementation source | Lane | Claim effect | Boundary |}\\
\noindent\texttt{|---|---|---|---|}\\
\noindent\texttt{| IMPL-LCR-TYPED-KERNEL | `normal\_form\_result` | Promotes L/C/R from notation to operational lane contract: LAdapter, CKernel, RChannel, strict policy gates, and lane classification. | This closes operational admission a}\\
\noindent\texttt{| IMPL-ANALOG-FORGE-WORKBENCH | `receipt\_bound\_internal\_result` | Converts the analog kit from a doctrine into a literal Python/PDF workbench: physical kit manifest, simulation layer, workbook sheets, receipts, and repor}\\

The immediate paper effect is stronger claim posture where these implementation rows satisfy the relevant lane. Remaining caution is reserved for specific claims that still lack an adapter, comparator, calibration target, or rerun receipt.

\subsection*{Appendix D. Resolved-State Closure Pass}

This pass removes false restrictions on the paper's claim posture. A row is no longer called open merely because a stronger future claim exists. The satisfied lane is closed now; only the stronger claim that lacks its required adapter, comparator, calibration, or falsifier remains as residue.

\subsubsection*{Closed Now}

\noindent\texttt{| Row | Lane | Resolved state | Remaining boundary |}\\
\noindent\texttt{|---|---|---|---|}\\
\noindent\texttt{| paper lift enumeration | `receipt\_bound\_internal\_result` | 8/8 LCR states succeeded; error walls=0; avg TarPit mass=0.506443. | This closes the paper-local finite lift readout, not every stronger physical interpretatio}\\
\noindent\texttt{| closure/lift governance | `normal\_form\_result` | The paper's object is closed as a replayable closure/lift state with explicit open-slot preservation. | The family closure is a structural lane; measured physics still r}\\
\noindent\texttt{| decade-1 tower | `receipt\_bound\_internal\_result` | The decade tower is resolved with avg TarPit mass=0.510370 and mass sum=40.829567. | The decade total is an internal tower metric, not a standalone universal physical }\\
\noindent\texttt{| family-01 cross-lift comparison | `cam\_crystal\_reapplication\_result` | Family tower binds FLCR-01, FLCR-11, FLCR-21, FLCR-31; avg TarPit mass=0.509490; error walls=0. | Cross-lift agreement strengthens routing and comp}\\
\noindent\texttt{| claim governance | `normal\_form\_result` | Claim-lane governance is resolved as the series' promotion rule. | Strong downstream physical claims still need their destination evidence lane. |}\\
\noindent\texttt{| analog access doctrine | `normal\_form\_result` | Analog workbooks are promoted from optional explanation to universal STEM access protocol. | Individual kits still require paper-local replay examples. |}\\

\subsubsection*{Claim Posture After This Pass}

`FLCR-01` should now be read as a resolved-state contribution for `Grounding Contract And Axiom Discipline` at its declared lane. The paper may state the strongest claim supported by these rows directly. It should reserve caution only for a specifically named stronger claim whose required evidence is absent.

\subsection*{8. Dependencies And Downstream Use}

This paper may be imported by later FLCR papers only through the claim lanes
above. The default downstream consumers are:

\noindent{}-- `FLCR-11` through `FLCR-18` for window, receipt, admission, CA, algebraic,\\
  curvature, residue, and exceptional machinery.
\noindent{}-- `FLCR-28` through `FLCR-30` for CAM/crystal, corpus ribbon, and universal\\
  normal-form intake.
\noindent{}-- `FLCR-31` through `FLCR-40` only after the LCR-native result has been cited\\
  and the translation claim has its own evidence lane.

\subsection*{9. Limitations And Falsifiers}

\noindent{}-- This paper does not itself prove later physics-facing claims; it defines how those claims are staged for later literal-physics adapters.\\
\noindent{}-- A high CAM score does not replace the need for attached receipts in the final publication package.\\
\noindent{}-- Imported mathematics must be cited as imported mathematics, not rebranded as internally derived.\\

A falsifier for this paper is any example that satisfies the declared input
conditions while violating the stated construction, receipt, or lane boundary.
An interpretation that merely wants a stronger downstream conclusion is not a
falsifier; it is an obligation routed to a later paper.

\subsection*{10. Publication Readiness Tasks}

\noindent{}-- Validator identities and receipt hashes are recorded in the manifest or receipt appendix.\\
\noindent{}-- External theorems require final citation entries before journal submission.\\
\noindent{}-- Every theorem, proposition, and protocol must retain a claim lane and evidence boundary.\\
\noindent{}-- The Markdown, LaTeX, and PDF forms are generated from the same canonical source.\\

\subsection*{Dual Positioning: Story And Formal Carrier}

This paper must be read in two synchronized ways. Sequentially, it explains why
the next state exists in the corpus story. Formally, it binds that state to
accepted mathematics, IRL formalism, code receipts, validators, CAM/crystal
lookups, and claim-lane boundaries.

Assigned original ribbon role(s): `00`/closure\_lift.

\noindent\texttt{| Original slot | Family | Lift | Role | Proof form |}\\
\noindent\texttt{|---|---|---:|---|---|}\\
\noindent\texttt{| `00` | closure\_lift | 0 | 0x ten-window closure/lift | aggregate receipt and open-slot preservation |}\\

The formal obligation is to state the strongest valid claim available for this
paper without letting interpretation silently change the addressed object. Any
author interpretation belongs beside the formalism as a declared relabeling,
bridge, or analog, and must be accompanied by the evidence lane that makes the
claim consumable by later papers.

\subsection*{State-Bound Proof Contract}

This paper receives: state emitted by prior slot 00 and reopened at original slot 00 (Established Grounding and Axiom Contract).

It must produce: formal result of original slot 00 plus the same-family lift path toward slot 10.

Semantic transition: bind the preceding window into a replayable state and expose the next lift without erasing residue.

Accepted formalism targets to bind in the journal rewrite:

\noindent{}-- conservative extension discipline\\
\noindent{}-- event sourcing and replay logs\\
\noindent{}-- finite receipt verification\\
\noindent{}-- dependency graph invariants\\

\noindent\texttt{| Slot | Family | Lift | Produced proof form |}\\
\noindent\texttt{|---|---|---:|---|}\\
\noindent\texttt{| `00` | closure\_lift | 0 | aggregate receipt and open-slot preservation |}\\

The formal carrier uses these accepted tools while preserving interpretive
claims as explicit relabelings, correspondences, or bridges. Claim promotion is admissible only when a receipt, standard citation,
CAM/crystal reapplication, normal-form result, calibration datum, or falsifier
boundary is attached.

\subsection*{Provenance And Crystal Evidence Integration}

This section integrates prior work, CAM/crystal projections, NSIT tooling,
standards-window evidence, and routed supplements as provenance-bearing
evidence. Source material is not treated as manuscript text; it is bound into
the paper only through claim lanes, receipts, validators, normal forms,
calibration rows, or explicit falsifier boundaries.

\subsubsection*{Expansion Rule For This Slot}

Past work binds the window as a replayable receipt stack: source, operation, result, residue, dependency, and lift boundary.

\subsubsection*{Primary Evidence Inputs}

\noindent\texttt{| Source | Placement reason |}\\
\noindent\texttt{|---|---|}\\
\noindent\texttt{| `00\_Established\_Grounding\_and\_Axiom\_Contract.md` | primary assigned source for the paper's formal spine |}\\
\noindent\texttt{| `supplements/00\_INTERNAL\_REASONING\_SUPPLEMENT.md` | paper-local reasoning supplement contributing to definitions, proof sketch, and limitations |}\\

\subsubsection*{Secondary And Orbital Evidence Inputs}

\noindent\texttt{| Source | Placement reason |}\\
\noindent\texttt{|---|---|}\\
\noindent\texttt{| `supplements/RECEIPT\_VERIFIER\_CATALOG.md` | cross-cutting supplement; paper-relevant rows, receipts, and guard language are bound through evidence lanes |}\\
\noindent\texttt{| `supplements/LATTICE\_FORGE\_UNIFICATION\_SUPPLEMENT.md` | cross-cutting supplement; paper-relevant rows, receipts, and guard language are bound through evidence lanes |}\\
\noindent\texttt{| `supplements/INTERNAL\_REASONING\_PYRAMID\_INDEX.md` | cross-cutting supplement; paper-relevant rows, receipts, and guard language are bound through evidence lanes |}\\
\noindent\texttt{| `supplements/INTERNAL\_REASONING\_5P\_WINDOW\_SUPPLEMENT.md` | cross-cutting supplement; paper-relevant rows, receipts, and guard language are bound through evidence lanes |}\\
\noindent\texttt{| `supplements/INTERNAL\_REASONING\_7P\_WINDOW\_SUPPLEMENT.md` | cross-cutting supplement; paper-relevant rows, receipts, and guard language are bound through evidence lanes |}\\
\noindent\texttt{| `supplements/INTERNAL\_REASONING\_9P\_WINDOW\_SUPPLEMENT.md` | cross-cutting supplement; paper-relevant rows, receipts, and guard language are bound through evidence lanes |}\\
\noindent\texttt{| `CAM\_CLAIM\_100\_SCORE\_AUDIT.md` | series-wide audit/score/closure evidence; import paper-specific rows and leave global context outside the body |}\\
\noindent\texttt{| `NSIT\_TOOL\_CLOSURE\_MATRIX.md` | series-wide audit/score/closure evidence; import paper-specific rows and leave global context outside the body |}\\
\noindent\texttt{| `NSIT\_REASONED\_CLOSURE\_PASS.md` | series-wide audit/score/closure evidence; import paper-specific rows and leave global context outside the body |}\\
\noindent\texttt{| `ORBITAL\_DATA\_CROSS\_POPULATION\_AUDIT.md` | series-wide audit/score/closure evidence; import paper-specific rows and leave global context outside the body |}\\
\noindent\texttt{| `AGENT\_CRYSTAL\_WORK\_HARVEST.md` | series-wide audit/score/closure evidence; import paper-specific rows and leave global context outside the body |}\\
\noindent\texttt{| `PAPER\_REWORKS\_CRYSTAL\_PROJECTION.json` | series-wide audit/score/closure evidence; import paper-specific rows and leave global context outside the body |}\\
\noindent\texttt{| Additional source routes | Complete routing is maintained in the source-placement appendix |}\\

\subsubsection*{Crystal And Standards Evidence To Bind}

\noindent{}-- Paper Reworks crystal projection: 33 paper markdown files, 9 supplements, 5 queues, 6 lattice-forge unification artifacts, 140 faces, 140 vignettes, and 280 CAM nodes.\\
\noindent{}-- Standards-window intake: 192/192 corpus conformance verdicts PASS; 140/142 obligations have candidate routes; 2/4/8/16/32 window lattice is available for cross-paper reads.\\
\noindent{}-- Agent/crystal harvest: agent-generated memory, runtime standards starter, current runtime build, repo-harvest CAM, notebook-derived bundles, and downloaded crystal payloads are orbital evidence surfaces.\\
\noindent{}-- NSIT inventory baseline: 220 validators, 1786 receipt/data artifacts, 1137 formal-paper-like artifacts, 114 supplements, and 86 memory/CAM artifacts.\\

\subsubsection*{Paper-Specific CAM/Score Rows}

\noindent\texttt{| 00 | 96 | system-composition closed | grounding contract, imported theorem burden, T3-T7 foundation, claim taxonomy | citation polish and cross-paper burden table finalization |}\\

\subsubsection*{Paper-Specific NSIT Closure Rows}

No direct NSIT row matched the assigned legacy papers in the first-pass matrix; validator matching by object name remains a receipt-attachment requirement.

\subsubsection*{Source Integration Summary}

The legacy source and internal reasoning supplement are treated as source
evidence rather than manuscript prose. Their contributions enter this paper as
claim-lane entries, receipt or validator references, finite-domain statements,
and limitation boundaries. Speculative or cross-domain interpretations remain
separate from closed local results until an explicit adapter, citation,
normal-form derivation, calibration row, or falsifier is attached.
\subsection*{Claim Binding Queue}

This queue records the evidence discipline for claim strengthening. It separates claims that already have support from residues that remain in falsifier or open-obligation lanes.

\noindent\texttt{| Queue | Lane | Promote now | Bind next | Residue |}\\
\noindent\texttt{|---|---|---|---|---|}\\
\noindent\texttt{| Formal-Methods Receipt And Governance Closure | `normal\_form\_result` | Promote claim-envelope, receipt, lane, causal-graph, and conformance obligations as formal-methods artifacts when standards reports are attached. |}\\

\subsubsection*{Binding Discipline}

Each queue item is represented as a delta:

\begin{verbatim}
local claim at paper time
-> attached later source/tool/receipt/theorem
-> revised representation
-> invariant boundary and remaining residue
\end{verbatim}

This preserves the sequential story while allowing later tools, crystals, and
standard formalisms to return through the proper lane.

\subsection*{Claim Delta Ledger}

This ledger applies the paper's claim-binding queues as explicit back-application
deltas. It keeps the sequential paper story intact while allowing later
receipts, standard formalisms, CAM/crystal routes, and validators to sharpen the
claim.

\noindent\texttt{| Delta | Local claim at paper time | Later evidence | Revised representation | Boundary kept |}\\
\noindent\texttt{|---|---|---|---|---|}\\
\noindent\texttt{| Formal Methods Receipt Closure | This paper may treat receipts, claim lanes, dependency graphs, and replay contracts as formal objects, not mere editorial apparatus. | NSIT standards, claim envelopes, content-addressed}\\

The revised representation records the strongest claim supported by the current evidence package. Promotion into theorem or proposition language occurs only when the named evidence is attached in the manifest or workbook replay.

\subsection*{Source-Backed Evidence Summary}

Routed archive and mirror cards are treated as provenance summaries rather than
body text. They identify candidate receipts, validators, theorem registries, or
supplemental source routes. A card strengthens this paper only when its
contribution is bound to a claim lane, evidence object, and boundary.

\noindent\texttt{| Evidence card | Score | Publication contribution | Boundary |}\\
\noindent\texttt{|---|---:|---|---|}\\
\noindent\texttt{| CQE-paper-10.50: Paper 10.50 - T10 Master Receipt Claim Contract | 85 | Paper 10.50 lets later papers cite the master receipt honestly. It makes T10 a powerful audit object without converting open lifts into hidden pro}\\
\noindent\texttt{| CQE-paper-10: Paper 10 - T10 Master Receipt | 81 | Paper 10 proves the first stack-level receipt theorem in the CQECMPLX suite. Its object is not a new physical mechanism. Its object is the proof that Paper 00 and Pape}\\
\noindent\texttt{| FORMAL: Paper 21 — C-Form: MorphForge / PolyForge / MorphoniX Gluon | 77 | **C = the morphic Gluon** — the token/number/shape/glyph transport Gluon that generalizes the ribbon transport to arbitrary symbolic tokens. In}\\
\noindent\texttt{| FINAL\_FORMAL\_PAPER: Complete Formal Claims: The Folded Strand | 75 | We present the complete closed-form claim set of the CQE\_CMPLX corpus: - 33 papers = 33 folding operations on a self-complementary strand - 144 tools}\\
\noindent\texttt{| SUMMARY-V-32-Theorems-Registry: Summary Paper V — The 32 Theorems in One Table: Closed-Form Registry | 75 | This paper is the **complete closed-form registry of all 32 theorems** in the CQE\_CMPLX corpus. Each theorem i}\\
\noindent\texttt{| Additional evidence cards | 86 total | The complete card inventory is maintained in the archive evidence matrix. | Cards are source-discovery aids until bound to paper-local evidence. |}\\

\subsection*{Journal Apparatus}

\subsubsection*{Article Type And Intended Venue Fit}

This manuscript is written as a formal-methods and mathematical-systems paper
inside the series *Formalizing LCR, Unifying Standard Models*. Its intended
reader is a technical reviewer who needs claim boundaries, reproducibility
routes, and cross-paper dependencies to be visible without relying on private
context.

\subsubsection*{Structured Contribution Statement}

\noindent\texttt{| Item | Statement |}\\
\noindent\texttt{|---|---|}\\
\noindent\texttt{| Publication ID | `FLCR-01` |}\\
\noindent\texttt{| Legacy source slot(s) | `00` |}\\
\noindent\texttt{| Ribbon role | 0x ten-window closure/lift |}\\
\noindent\texttt{| Proof form | aggregate receipt and open-slot preservation |}\\
\noindent\texttt{| Received state | state emitted by prior slot 00 and reopened at original slot 00 (Established Grounding and Axiom Contract) |}\\
\noindent\texttt{| Produced state | formal result of original slot 00 plus the same-family lift path toward slot 10 |}\\

\subsubsection*{Claim-Evidence Table}

\noindent\texttt{| Claim | Lane | Evidence to attach | Boundary |}\\
\noindent\texttt{|---|---|---|---|}\\
\noindent\texttt{| Theorem 1.1 | `normal\_form\_result` | Receipt, source card, validator, citation, or CAM route named in the paper manifest | The FLCR corpus is admissible only as a typed extension stack: every promoted claim must declar}\\
\noindent\texttt{| Proposition 1.2 | `standard\_theorem\_citation\_bound\_result` | Receipt, source card, validator, citation, or CAM route named in the paper manifest | For finite domains such as the eight binary LCR states, exhaustive enum}\\
\noindent\texttt{| Protocol 1.3 | `falsifier\_or\_open\_obligation` | Receipt, source card, validator, citation, or CAM route named in the paper manifest | A claim without a receipt, citation, calibration datum, CAM route, or falsifier boun}\\

\subsubsection*{Figure Plan}

\noindent\texttt{| Figure | Purpose | Caption |}\\
\noindent\texttt{|---|---|---|}\\
\noindent\texttt{| FLCR-01-F1 | Receipt stack and lift boundary | Schematic showing how `FLCR-01` turns its received state into the produced state while preserving claim lanes and residue boundaries. |}\\
\noindent\texttt{| FLCR-01-F2 | Evidence routing map | Diagram of source papers, archive cards, CAM/crystal routes, validators, and workbook replay paths used by this manuscript. |}\\
\noindent\texttt{| FLCR-01-F3 | Same-family lift context | Diagram placing this paper in the original `00-79` ribbon family and showing predecessor/successor slots. |}\\

\subsubsection*{In-Text Figure FLCR-01-F1: Receipt stack and lift boundary}

\begin{verbatim}
received state
  -> Receipt stack and lift boundary
  -> formal claim lane
  -> evidence object / receipt / citation
  -> produced state
  -> remaining residue
\end{verbatim}

\subsubsection*{Table Plan}

\noindent\texttt{| Table | Purpose |}\\
\noindent\texttt{|---|---|}\\
\noindent\texttt{| FLCR-01-T1 | Receipt/lift table |}\\
\noindent\texttt{| FLCR-01-T2 | Claim-lane/evidence-boundary table |}\\
\noindent\texttt{| FLCR-01-T3 | Archive evidence card and source-backed upgrade table |}\\
\noindent\texttt{| FLCR-01-T4 | Workbook replay and falsifier table |}\\

\subsubsection*{Worked Example}

Trace one prior-window claim through source, operation, receipt, residue, and downstream import. The example must name the input object, the operation,
the evidence object, the allowed revised claim, and the remaining boundary.

\subsubsection*{Nomenclature And Glossary}

\noindent\texttt{| Term | Paper-local meaning |}\\
\noindent\texttt{|---|---|}\\
\noindent\texttt{| CAM | Defined by this paper as part of the `closure\_lift` proof role; final definition is recorded in the paper-local glossary. |}\\
\noindent\texttt{| claim lane | Defined by this paper as part of the `closure\_lift` proof role; final definition is recorded in the paper-local glossary. |}\\
\noindent\texttt{| crystal | Defined by this paper as part of the `closure\_lift` proof role; final definition is recorded in the paper-local glossary. |}\\
\noindent\texttt{| dependency graph | Defined by this paper as part of the `closure\_lift` proof role; final definition is recorded in the paper-local glossary. |}\\
\noindent\texttt{| lift boundary | Defined by this paper as part of the `closure\_lift` proof role; final definition is recorded in the paper-local glossary. |}\\
\noindent\texttt{| open-slot preservation | Defined by this paper as part of the `closure\_lift` proof role; final definition is recorded in the paper-local glossary. |}\\
\noindent\texttt{| receipt | Defined by this paper as part of the `closure\_lift` proof role; final definition is recorded in the paper-local glossary. |}\\

\subsubsection*{Appendix FLCR-01-A: Evidence Package}

The appendix records:

1. Legacy source excerpts or source-card references used in the final claim ledger.
2. Receipt and verifier paths, including pass/fail/partial status.
3. CAM/crystal query or projection rows used by the paper, with source identity and boundary.
4. Standards-window and NSIT evidence used for conformance or routing.
5. Falsifier, calibration, or external-data rows that remain unresolved.

\subsubsection*{Data And Code Availability}

The publication package contains the canonical Markdown sources, manifests, source-routing matrices, validation reports, and generated PDF/TeX outputs.

Code and verifier references are cited through manifests, archive evidence cards, receipt catalogs, or workbook replay tables. External datasets or
standards citations must be attached before any calibration-bound claim is
promoted.

\subsubsection*{Reviewer Checklist}

\noindent\texttt{| Check | Reviewer question |}\\
\noindent\texttt{|---|---|}\\
\noindent\texttt{| Claim lane | Does every strong claim declare its lane and evidence type? |}\\
\noindent\texttt{| Boundary | Does the paper preserve what it does not prove? |}\\
\noindent\texttt{| Reproducibility | Is there a receipt, verifier, source card, citation, or replay table for the claim? |}\\
\noindent\texttt{| Cross-paper import | Does the paper cite the prior FLCR/OPR slot before consuming its result? |}\\
\noindent\texttt{| Interpretation | Does the analog layer preserve the addressed state while changing labels/access paths? |}\\

\subsubsection*{Declarations}

No human-subjects, clinical, or animal-research protocol is asserted by this
paper. External empirical claims require separate dataset, calibration, or
domain-validation attachments. Competing-interest and funding statements are recorded in the submission metadata.

\subsection*{11. Publication Integration Addendum}

\subsubsection*{11.1 Role In The Series}

`FLCR-01` (`Grounding Contract And Axiom Discipline`) occupies the `closure\_lift` role at lift depth `0`.
The paper receives state emitted by prior slot 00 and reopened at original slot 00 (Established Grounding and Axiom Contract). It produces formal result of original slot 00 plus the same-family lift path toward slot 10. The state transition is:
bind the preceding window into a replayable state and expose the next lift without erasing residue.

The paper-local proof form is:

\begin{verbatim}
aggregate receipt and open-slot preservation
\end{verbatim}

The integration result is a receipt-bound closure state with explicit open-slot preservation. This addendum records how that result is
consumed by the publication series without relying on editorial instructions or
private working context.

\subsubsection*{11.2 Formal Carrier}

\noindent\texttt{| Formal carrier | Function |}\\
\noindent\texttt{|---|---|}\\
\noindent\texttt{| conservative extension discipline | Formal carrier for the paper-local result. |}\\
\noindent\texttt{| event sourcing and replay logs | Formal carrier for the paper-local result. |}\\
\noindent\texttt{| finite receipt verification | Formal carrier for the paper-local result. |}\\
\noindent\texttt{| dependency graph invariants | Formal carrier for the paper-local result. |}\\

The LCR, CAM, crystal, forge, and analog vocabulary functions as a typed access
layer over these carriers. A relabeling is admissible only when the addressed
object, evidence lane, boundary, and downstream import rule remain attached.

\subsubsection*{11.3 Claim Ledger}

\noindent\texttt{| Claim | Lane | Statement |}\\
\noindent\texttt{|---|---|---|}\\
\noindent\texttt{| Theorem 1.1 | `normal\_form\_result` | The FLCR corpus is admissible only as a typed extension stack: every promoted claim must declare its lane, source, evidence, and boundary. |}\\
\noindent\texttt{| Proposition 1.2 | `standard\_theorem\_citation\_bound\_result` | For finite domains such as the eight binary LCR states, exhaustive enumeration is complete proof of the enumerated property. |}\\
\noindent\texttt{| Protocol 1.3 | `falsifier\_or\_open\_obligation` | A claim without a receipt, citation, calibration datum, CAM route, or falsifier boundary remains a staged obligation. |}\\

These claims are consumed at their stated lane. A stronger downstream statement
creates a new claim envelope with its own evidence object.

\subsubsection*{11.4 Evidence Package}

\noindent\texttt{| Evidence class | Routed items | Status |}\\
\noindent\texttt{|---|---|---|}\\
\noindent\texttt{| Legacy sources | `00\_Established\_Grounding\_and\_Axiom\_Contract.md`, `supplements/00\_INTERNAL\_REASONING\_SUPPLEMENT.md` | routed evidence |}\\
\noindent\texttt{| Evidence inputs | `CAM\_CLAIM\_100\_SCORE\_AUDIT.md`, `NSIT\_TOOL\_CLOSURE\_MATRIX.md`, `NSIT\_REASONED\_CLOSURE\_PASS.md`, `PAPER\_REWORKS\_CRYSTAL\_PROJECTION.json` | routed evidence |}\\
\noindent\texttt{| CAM/crystal sources | `PAPER\_REWORKS\_CRYSTAL\_PROJECTION.json`, `AGENT\_CRYSTAL\_WORK\_HARVEST.json`, `runtime standards artifact`, `notebook-derived source bundle` | routed evidence |}\\
\noindent\texttt{| Standards-window inputs | `window\_summary.json`, `node DBs`, `window DBs` | routed evidence |}\\
\noindent\texttt{| Receipts | external requirement | not attached in this source package |}\\
\noindent\texttt{| Validators | external requirement | not attached in this source package |}\\

Standards-window, runtime, and notebook-derived artifacts are treated
as evidence-routing surfaces. They support source discovery, conformance
checking, and reapplication, but they do not replace citations, receipts,
normal-form derivations, calibration rows, or falsifiers.

\subsubsection*{11.5 Results Representation}

The paper's results section is organized around five publication objects:

1. the exact object acted on at this slot;
2. the admissible operation or transformation;
3. the receipt, enumeration, derivation, lookup, or external theorem supporting
   the operation;
4. the residue or boundary left after the result;
5. the downstream import rule.

\subsubsection*{11.6 Figures And Tables}

\noindent\texttt{| Figure | Role |}\\
\noindent\texttt{|---|---|}\\
\noindent\texttt{| FLCR-01-F1: Receipt stack and lift boundary | Visualizes the proof object, routing, or boundary. |}\\
\noindent\texttt{| FLCR-01-F2: Evidence routing map | Visualizes the proof object, routing, or boundary. |}\\
\noindent\texttt{| FLCR-01-F3: Same-family lift context | Visualizes the proof object, routing, or boundary. |}\\

\noindent\texttt{| Table | Role |}\\
\noindent\texttt{|---|---|}\\
\noindent\texttt{| FLCR-01-T1: Receipt/lift table | Records claim lanes, evidence, sources, or falsifiers. |}\\
\noindent\texttt{| FLCR-01-T2: Claim-lane/evidence-boundary table | Records claim lanes, evidence, sources, or falsifiers. |}\\
\noindent\texttt{| FLCR-01-T3: Archive evidence card and source-backed upgrade table | Records claim lanes, evidence, sources, or falsifiers. |}\\
\noindent\texttt{| FLCR-01-T4: Workbook replay and falsifier table | Records claim lanes, evidence, sources, or falsifiers. |}\\

\subsubsection*{11.7 Limits And Falsifiers}

The paper proves only the object and domain declared in its claim ledger.
Physical, Standard Model, observer, calibration, or synthesis claims require
the additional lane evidence stated by their destination papers. CAM/crystal
and forge language is operational only when represented as an address, lookup,
receipt, validator, or reapplication path.
\end{document}
